% ===========================================================================
%  Phase Field Theory of Inequality and Consensus
%  不平等与共识的相场理论
% ===========================================================================
%  Core: Phase field methods model the evolution of interfaces between
%  different "phases" of the social/audit system — honest regions (g=0)
%  and dishonest regions (g≠0) form domains separated by diffuse interfaces.
%  相场方法建模社会/审计系统中不同"相"之间界面的演化——
%  诚实地带(g=0)和不诚实地带(g≠0)形成由扩散界面分隔的畴结构。
% ===========================================================================

\documentclass[11pt,a4paper]{article}

% ===========================================================================
%  Packages
% ===========================================================================
\usepackage[utf8]{inputenc}
\usepackage[T1]{fontenc}
\usepackage{amsmath,amssymb,amsthm}
\usepackage{mathtools}
\usepackage{mathrsfs}
\usepackage{bbm}
\usepackage{bm}
\usepackage{graphicx}
\usepackage{xcolor}
\usepackage{hyperref}
\usepackage{enumitem}
\usepackage{booktabs}
\usepackage{tikz}
\usepackage{tikz-cd}
\usepackage{float}
\usepackage{caption}
\usepackage{subcaption}
\usepackage{geometry}
\usepackage[colorinlistoftodos]{todonotes}
\usepackage{xeCJK}

% Page geometry
\geometry{margin=1.2in,top=1in,bottom=1in}

% ===========================================================================
%  Theorem environments
% ===========================================================================
\theoremstyle{plain}
\newtheorem{theorem}{Theorem}[section]
\newtheorem{lemma}[theorem]{Lemma}
\newtheorem{proposition}[theorem]{Proposition}
\newtheorem{corollary}[theorem]{Corollary}
\newtheorem{conjecture}[theorem]{Conjecture}

\theoremstyle{definition}
\newtheorem{definition}[theorem]{Definition}
\newtheorem{example}[theorem]{Example}
\newtheorem{remark}[theorem]{Remark}
\newtheorem{protocol}[theorem]{Protocol}

% ===========================================================================
%  Math operators and notation
% ===========================================================================
\DeclareMathOperator{\im}{im}
\DeclareMathOperator{\vol}{vol}
\DeclareMathOperator{\diam}{diam}
\DeclareMathOperator{\supp}{supp}
\DeclareMathOperator{\grad}{grad}
\DeclareMathOperator{\divg}{div}
\DeclareMathOperator{\lap}{\Delta}
\DeclareMathOperator{\Yajie}{Yajie}
\DeclareMathOperator{\Cercis}{Cercis}

\newcommand{\R}{\mathbb{R}}
\newcommand{\C}{\mathbb{C}}
\newcommand{\Z}{\mathbb{Z}}
\newcommand{\N}{\mathbb{N}}
\newcommand{\Q}{\mathbb{Q}}
\newcommand{\cM}{\mathcal{M}}
\newcommand{\cS}{\mathcal{S}}
\newcommand{\cF}{\mathcal{F}}
\newcommand{\cA}{\mathcal{A}}
\newcommand{\cI}{\mathcal{I}}
\newcommand{\cD}{\mathcal{D}}
\newcommand{\cP}{\mathcal{P}}
\newcommand{\cH}{\mathcal{H}}
\newcommand{\cC}{\mathcal{C}}
\newcommand{\cL}{\mathcal{L}}
\newcommand{\cO}{\mathcal{O}}
\newcommand{\cG}{\mathcal{G}}
\newcommand{\cB}{\mathcal{B}}
\newcommand{\cV}{\mathcal{V}}
\newcommand{\cK}{\mathcal{K}}
\newcommand{\cE}{\mathcal{E}}

% Phase-related notation
\newcommand{\gf}{\mathbf{g}}          % attitude/gauge field
\newcommand{\Sf}{S}                   % potential field
\newcommand{\deltac}{\delta_c}         % critical thickness
\newcommand{\phasesep}{\mathcal{I}}    % phase separator
\newcommand{\Rc}{R_c}                  % critical radius

% ===========================================================================
%  Title and Author
% ===========================================================================
\title{
  {\large\textbf{Phase Field Theory of Inequality and Consensus}}\\[4pt]
  {\large\textbf{Allen-Cahn / Cahn-Hilliard Dynamics on SCX Potential Manifolds}}\\[8pt]
  {\normalsize\textit{不平等与共识的相场理论}}\\[4pt]
  {\normalsize\textit{SCX势能流形上的Allen-Cahn/Cahn-Hilliard动力学}}\\[4pt]
  {\footnotesize Phase Separation, Nucleation, Coarsening, and Theorem Connections}\\[2pt]
  {\footnotesize\textit{相分离、成核、粗化与定理联系}}
}

\author{
  SCX Research Group\\[2pt]
  \texttt{Xiaogan Supercomputing Center}\\[4pt]
  \small\textit{Version 1.0 — July 2026}
}
\date{}

% ===========================================================================
\begin{document}
\maketitle

% ===========================================================================
%  Abstract
% ===========================================================================
\begin{abstract}
\noindent
We formulate the dynamics of inequality and consensus in the SCX (Situs Consensus
eXpert) framework as a \textbf{phase field theory}. The attitude/gauge field
$\gf(x,t)$ and the potential field $\Sf(x,t)$ are treated as coupled order
parameters on the Situs manifold $\cM$. Drawing from the classical
Allen-Cahn and Cahn-Hilliard equations of materials science, we derive:

\begin{enumerate}[label=(\arabic*)]
  \item \textbf{Allen-Cahn dynamics for $\gf$:} $\partial_t \gf = -\delta\cF/\delta\gf$
        — the gauge field evolves via non-conserved gradient flow to minimize a
        Ginzburg-Landau free energy, forming honest ($\gf=\mathbf{0}$) and
        dishonest ($\gf \neq \mathbf{0}$) domains separated by diffuse interfaces
        of thickness $\deltac$;

  \item \textbf{Cahn-Hilliard dynamics for $\Sf$:} $\partial_t \Sf = \nabla \cdot
        (M(\Sf) \nabla (\delta\cF/\delta\Sf))$ — the potential field undergoes
        mass-conserving evolution, redistributing potential while preserving total
        ``social energy,'' directly modeling inequality dynamics;

  \item \textbf{Three-phase diagram} in $(\|\gf\|, \Sf)$ space: (a) EQUALITY phase
        — $\|\gf\| \approx 0$, smooth $\Sf$; (b) STRATIFIED phase — $\|\gf\| > 0$,
        $\Sf$ exhibits plateaus; (c) EXPLOSIVE phase — $\|\gf\|$ diverges, $\Sf$
        gradients blow up. Phase boundaries are the \textbf{inequality phase
        transition lines};

  \item \textbf{Nucleation and growth:} A small honest cluster (``audit nucleus,''
        $\gf = \mathbf{0}$ domain) can grow only if its radius exceeds the critical
        radius $\Rc \propto \sigma / \Delta f$, where $\sigma$ is the interface
        energy and $\Delta f$ is the bulk free energy gain. We derive $\Rc$ from
        the competition between surface and volume terms;

  \item \textbf{Coarsening dynamics:} After a social ``quench'' (rapid institutional
        change), honest domains coarsen via Ostwald ripening — large honest domains
        grow at the expense of small ones. The characteristic domain size follows
        $L(t) \sim t^{1/3}$, the classical Lifshitz-Slyozov-Wagner (LSW) scaling;

  \item \textbf{Theorem connections:} Theorem~11 (attitude singularity) corresponds
        to a first-order phase transition at the spinodal; Theorem~10 (boundary
        locking) maps to \textbf{pinning of domain walls} at heterogeneities;
        Theorem~12 (Matthew staircase) maps to \textbf{nucleation on defects} —
        staircase steps act as preferential nucleation sites that accelerate
        system destabilization.
\end{enumerate}

We provide a complete verification script (\texttt{verify\_phase\_field.py})
that numerically solves the coupled Allen-Cahn / Cahn-Hilliard system on a
2D grid, confirms the $t^{1/3}$ coarsening scaling, computes the phase
diagram, and validates the critical nucleus prediction.

\bigskip
\noindent
\textit{我们将SCX（Situs共识专家）框架中的不平等与共识动力学表述为\textbf{相场理论}。态度/规范场 $\gf(x,t)$ 和势能场 $\Sf(x,t)$ 被处理为Situs流形 $\cM$ 上的耦合序参量。借鉴材料科学中经典的Allen-Cahn和Cahn-Hilliard方程，我们推导出：(1) $\gf$ 的\textbf{Allen-Cahn动力学} — 规范场通过非守恒梯度流演化以最小化Ginzburg-Landau自由能，形成诚实地带($\gf=\mathbf{0}$)和不诚实地带($\gf\neq\mathbf{0}$)的畴结构；(2) $\Sf$ 的\textbf{Cahn-Hilliard动力学} — 势能场经历质量守恒的演化，在保持总"社会能量"的同时重新分配势能；(3) 在 $(\|\gf\|, \Sf)$ 空间中的\textbf{三相图} — 平等相、分层相和爆炸相；(4) \textbf{成核与生长} — 审计核的临界半径；(5) \textbf{粗化动力学} — $L(t) \sim t^{1/3}$ 的LSW标度律；(6) \textbf{定理联系} — Thm11对应相变、Thm10对应畴壁钉扎、Thm12对应缺陷成核。}

\bigskip
\noindent
\textbf{Keywords:} phase field theory, Allen-Cahn equation, Cahn-Hilliard equation,
inequality dynamics, consensus, nucleation, coarsening, Ostwald ripening,
spinodal decomposition, SCX framework

\noindent
\textbf{关键词：} 相场理论，Allen-Cahn方程，Cahn-Hilliard方程，不平等动力学，
共识，成核，粗化，Ostwald熟化，旋节分解，SCX框架
\end{abstract}

% ===========================================================================
\tableofcontents

% ===========================================================================
%  SECTION 1: INTRODUCTION
% ===========================================================================
\section{Introduction: Inequality as a Phase Ordering Problem}
\section*{引言：不平等作为相有序化问题}

% ---------------------------------------------------------------------------
\subsection{Motivation from Materials Science}
\subsection*{来自材料科学的动机}

In materials science, \textbf{phase field methods} provide a unified framework
for describing the spatiotemporal evolution of microstructures — the spatial
arrangement of different phases (solid/liquid, ordered/disordered, grain
orientations) that determines macroscopic material properties
\cite{chen2002, provatas2011}. The key insight is to replace sharp interfaces
with \textbf{diffuse interfaces} described by continuous order parameter fields,
and to let these fields evolve according to thermodynamic driving forces encoded
in a free energy functional.

The two foundational equations are:

\begin{enumerate}[label=(\arabic*)]
  \item \textbf{Allen-Cahn equation} (non-conserved order parameter):
  \begin{equation}
    \frac{\partial \phi}{\partial t} = -M_\phi \frac{\delta \cF}{\delta \phi}
    \label{eq:ac_classic}
  \end{equation}
  describes the evolution of a non-conserved field $\phi$ (e.g., crystal
  orientation, magnetic spin) toward free energy minima. The dynamics produce
  domain growth via curvature-driven interfacial motion — boundaries move
  to reduce total interface area.

  \item \textbf{Cahn-Hilliard equation} (conserved order parameter):
  \begin{equation}
    \frac{\partial c}{\partial t} = \nabla \cdot \left(M_c(c) \nabla
    \frac{\delta \cF}{\delta c}\right)
    \label{eq:ch_classic}
  \end{equation}
  describes the evolution of a conserved field $c$ (e.g., concentration)
  where $\int c \, dx$ is constant. The dynamics produce spinodal decomposition
  and coarsening — the system phase-separates into domains that grow via
  diffusion-limited ripening.
\end{enumerate}

\bigskip
\noindent
\textit{在材料科学中，\textbf{相场方法}提供了描述微结构时空演化的统一框架。核心思想是用连续序参量场描述的\textbf{扩散界面}替代尖锐界面，并让这些场根据自由能泛函编码的热力学驱动力演化。两个基础方程是：(1) \textbf{Allen-Cahn方程} — 描述非守恒场向自由能极小值的演化，产生曲率驱动的界面运动；(2) \textbf{Cahn-Hilliard方程} — 描述守恒场的演化，产生旋节分解和粗化。}

% ---------------------------------------------------------------------------
\subsection{The SCX Analogy: Two Coupled Fields}
\subsection*{SCX类比：两个耦合场}

In the SCX framework \cite{scx_framework, scx_moe_gauge}, the social/audit
system is characterized by two fundamental fields on the Situs manifold $\cM$:

\begin{definition}[SCX Phase Fields / SCX相场]
\label{def:phase_fields}
\begin{enumerate}[label=(\roman*)]
  \item \textbf{Gauge field (态度场)} $\gf: \cM \times [0,\infty) \to \R^{d_g}$ —
  encodes the ``coordinate system bias'' of each agent/expert. Honest agents
  have $\gf = \mathbf{0}$ (their assessments align with the consensus origin);
  dishonest or biased agents have $\gf \neq \mathbf{0}$. The norm $\|\gf\|$
  measures the \textbf{degree of deviation} from the gauge-fixed consensus.

  \item \textbf{Potential field (势能场)} $\Sf: \cM \times [0,\infty) \to \R_{\ge 0}$ —
  encodes the accumulated ``stored social energy'' — wealth, status, resources,
  information — at each point. $\Sf$ is a \textbf{conserved quantity} in the
  sense that total potential $\int_\cM \Sf(x,t)\,dx$ changes only through
  external injection or extraction, not through internal redistribution alone.
\end{enumerate}
\end{definition}

The gauge fixing condition of SCX — $\sum_m \gf_m = \mathbf{0}$ — is exactly
the condition that there is no net bias in the system. When this condition is
violated locally, $\gf$-domains form: clusters of agents whose coordinate
systems are collectively shifted in the same direction. These $\gf$-domains
are separated from honest regions by \textbf{gauge domain walls} — interfaces
across which the bias field transitions from $\mathbf{0}$ to some non-zero
value.

\bigskip
\noindent
\textit{在SCX框架中，社会/审计系统由两个基本场刻画：(i) \textbf{规范场} $\gf$ — 编码每个智能体的"坐标系偏差"；(ii) \textbf{势能场} $\Sf$ — 编码累积的"社会能量"。SCX的规范固定条件 $\sum_m \gf_m = \mathbf{0}$ 正是系统无净偏差的条件。当此条件被局部违反时，$\gf$-畴形成——由\textbf{规范畴壁}分隔的偏差集群。}

% ---------------------------------------------------------------------------
\subsection{From Static Geometry to Dynamic Evolution}
\subsection*{从静态几何到动态演化}

Previous SCX papers have focused on the \textbf{static geometry} of the potential
surface: the existence of singularities (Thm~11), the instability of boundaries
(Thm~10), and the staircase structure of inequality (Thm~12). This paper takes
the next step: \textbf{dynamics}. How do $\gf$-domains form, grow, interact,
and coarsen? What drives the system from the EQUALITY phase to the STRATIFIED
phase, and from there to the EXPLOSIVE phase? Can we predict the critical
conditions for phase transitions?

The phase field formalism answers these questions by providing:

\begin{enumerate}[label=(\arabic*)]
  \item A \textbf{free energy functional} $\cF[\gf, \Sf]$ that encodes the
        thermodynamic cost of inhomogeneity and inequality;
  \item \textbf{Equations of motion} derived from variational principles;
  \item A \textbf{phase diagram} mapping stability regions in the
        $(\|\gf\|, \Sf)$ plane;
  \item \textbf{Kinetic pathways} — nucleation, growth, coarsening — that
        describe how the system traverses the phase diagram.
\end{enumerate}

\bigskip
\noindent
\textit{先前SCX论文聚焦于势能面的\textbf{静态几何}：奇点存在性(Thm11)、边界不稳定性(Thm10)、不平等的阶梯结构(Thm12)。本文迈出下一步：\textbf{动力学}。$\gf$-畴如何形成、生长、相互作用和粗化？系统如何从平等相到分层相再到爆炸相？相场形式化通过提供自由能泛函、运动方程、相图和动力学路径来回答这些问题。}

% ===========================================================================
%  SECTION 2: FREE ENERGY FUNCTIONAL
% ===========================================================================
\section{Free Energy Functional of the SCX System}
\section{SCX系统的自由能泛函}

% ---------------------------------------------------------------------------
\subsection{The Ginzburg-Landau Free Energy}
\subsection*{Ginzburg-Landau自由能}

We postulate that the SCX system minimizes a free energy functional of the
Ginzburg-Landau type, with separate contributions from the gauge field and
the potential field, plus a coupling term:

\begin{definition}[SCX Free Energy Functional / SCX自由能泛函]
\label{def:free_energy}
For a Situs manifold $\cM \subset \R^d$, the total free energy is:

\begin{equation}
  \cF[\gf, \Sf] = \int_\cM \left[
    f_g(\gf) + \frac{\kappa_g}{2} |\nabla \gf|^2
    + f_S(\Sf) + \frac{\kappa_S}{2} |\nabla \Sf|^2
    + \lambda \, \Gamma(\gf, \Sf)
  \right] dx
  \label{eq:total_free_energy}
\end{equation}

where:
\begin{itemize}
  \item $f_g(\gf) = \frac{A}{4} \|\gf\|^4 - \frac{B}{2} \|\gf\|^2$ —
        the \textbf{gauge double-well potential}. $A, B > 0$ are material
        parameters. The minima are at $\|\gf\| = 0$ (honest phase) and
        $\|\gf\| = \sqrt{B/A}$ (dishonest phase). The barrier height is
        $B^2/(4A)$.

  \item $\frac{\kappa_g}{2} |\nabla \gf|^2$ — the \textbf{gauge gradient
        energy}, penalizing sharp spatial variations in bias. $\kappa_g > 0$
        is the gauge stiffness. This term gives domain walls a finite thickness.

  \item $f_S(\Sf) = \frac{\alpha}{2} (\Sf - \Sf_0)^2 + \frac{\beta}{4} \Sf^4$ —
        the \textbf{potential energy density}. $\Sf_0$ is the ``natural''
        potential level. The quartic term prevents runaway growth.

  \item $\frac{\kappa_S}{2} |\nabla \Sf|^2$ — the \textbf{potential gradient
        energy}, penalizing sharp potential gradients (i.e., high inequality).
        $\kappa_S > 0$ is the social mobility stiffness.

  \item $\lambda \, \Gamma(\gf, \Sf)$ — the \textbf{gauge-potential coupling}.
        $\lambda \ge 0$ is the coupling strength. We propose:
        \begin{equation}
          \Gamma(\gf, \Sf) = \|\gf\|^2 \cdot \Sf
          \label{eq:coupling}
        \end{equation}
        This encodes the observation that high bias ($\|\gf\|$) combined with
        high potential ($\Sf$) is the most dangerous configuration — the
        ``double explosion'' of Theorem~11.
\end{itemize}
\end{definition}

\bigskip
\noindent
\textit{SCX自由能泛函包含：规范双阱势 $f_g$（极值在诚实地带 $\|\gf\|=0$ 和不诚实地带 $\|\gf\|=\sqrt{B/A}$）；规范梯度能（$\kappa_g$ 项）为畴壁赋予有限厚度；势能密度 $f_S$；势能梯度能（$\kappa_S$ 项）惩罚高不平等；以及规范-势能耦合项 $\lambda \|\gf\|^2 \Sf$ — 编码Thm11的"高偏差+高势能=双重爆炸"。}

% ---------------------------------------------------------------------------
\subsection{Interpretation of Terms}
\subsection*{各项的物理解释}

\begin{enumerate}[label=(\arabic*)]
  \item \textbf{Gauge double-well:} $f_g$ has two minima, reflecting the
        bistable nature of honesty/dishonesty. An agent's bias can be driven
        to either state by local social pressure. The barrier between them
        represents the ``activation energy'' for changing one's ethical stance.

  \item \textbf{Gauge gradient energy:} The term $\frac{\kappa_g}{2}|\nabla\gf|^2$
        penalizes the \textbf{interface} between honest and dishonest domains.
        In social terms, this is the ``cognitive dissonance cost'' of being
        at a boundary — agents whose neighbors have different biases
        experience tension. The parameter $\kappa_g$ measures social cohesion:
        large $\kappa_g$ means thick, diffuse boundaries (tolerant society);
        small $\kappa_g$ means sharp boundaries (polarized society).

  \item \textbf{Potential free energy:} $f_S$ encodes that there is an
        ``optimal'' potential level $\Sf_0$, with quadratic penalties for
        deviation. The quartic term $\frac{\beta}{4}\Sf^4$ ensures that
        extreme wealth concentration has super-linear social cost.

  \item \textbf{Coupling:} $\lambda \|\gf\|^2 \Sf$ is the crucial term that
        links bias and inequality. High bias \textit{and} high potential is
        explosive; high bias with low potential is less dangerous; low bias
        with high potential is benign. This is the phase-field encoding of
        the ``attitude-potential double explosion'' from Theorem~11.
\end{enumerate}

\bigskip
\noindent
\textit{规范双阱势反映诚实/不诚实的双稳态。规范梯度能编码"认知失调成本"——$\kappa_g$ 大意味着容忍社会（厚扩散边界），$\kappa_g$ 小意味着极化社会（尖锐边界）。势能自由能确保极端财富集中具有超线性的社会成本。耦合项 $\lambda \|\gf\|^2 \Sf$ 是连接偏差与不平等的关键——高偏差+高势能=爆炸。}

% ===========================================================================
%  SECTION 3: ALLEN-CAHN DYNAMICS FOR THE GAUGE FIELD
% ===========================================================================
\section{Allen-Cahn Dynamics for the Gauge Field}
\section{规范场的Allen-Cahn动力学}

% ---------------------------------------------------------------------------
\subsection{Equation of Motion}
\subsection*{运动方程}

The gauge field $\gf$ is a \textbf{non-conserved} order parameter — there is
no conservation law for total ``bias'' in the system. Honest agents can become
dishonest and vice versa without any global constraint. Therefore $\gf$ evolves
via \textbf{Allen-Cahn (Model A)} dynamics — gradient descent on the free
energy:

\begin{definition}[Gauge Allen-Cahn Equation / 规范Allen-Cahn方程]
\label{def:gauge_ac}
\begin{equation}
  \frac{\partial \gf}{\partial t} = -M_g \frac{\delta \cF}{\delta \gf}
  = -M_g \left[
    \left(A \|\gf\|^2 - B\right) \gf
    - \kappa_g \lap \gf
    + 2\lambda \Sf \, \gf
  \right]
  \label{eq:gauge_ac}
\end{equation}
where $M_g > 0$ is the gauge mobility (rate at which agents adjust their bias).
\end{definition}

The variational derivative is computed component-wise:
\begin{equation}
  \frac{\delta \cF}{\delta \gf} = \frac{\partial f_g}{\partial \gf}
  - \kappa_g \lap \gf + \lambda \frac{\partial \Gamma}{\partial \gf}
  = (A\|\gf\|^2 - B)\gf - \kappa_g \lap \gf + 2\lambda \Sf \gf
  \label{eq:variational_g}
\end{equation}

\bigskip
\noindent
\textit{规范场 $\gf$ 是非守恒序参量——系统中"偏差"总量没有守恒律。因此 $\gf$ 通过Allen-Cahn (Model A) 动力学演化：梯度下降于自由能。运动方程包含三个驱动力：双阱势的体积力 $(A\|\gf\|^2-B)\gf$、扩散项 $\kappa_g\lap\gf$、以及耦合项 $2\lambda\Sf\gf$。}

% ---------------------------------------------------------------------------
\subsection{Interface Structure: The Gauge Domain Wall}
\subsection*{界面结构：规范畴壁}

Consider a one-dimensional stationary interface between an honest domain
($\gf=0$ at $x \to -\infty$) and a dishonest domain ($\|\gf\|=\sqrt{B/A}$
at $x \to +\infty$), with $\Sf$ held constant. The equilibrium profile
satisfies:

\begin{equation}
  \kappa_g \frac{d^2 \gf}{dx^2} = \frac{\partial f_g^{\text{eff}}}{\partial \gf}
  \label{eq:equilibrium_wall}
\end{equation}

where $f_g^{\text{eff}}(\gf) = f_g(\gf) + \lambda \|\gf\|^2 \Sf$. In the
decoupled limit ($\lambda \to 0$), the solution is the classical
Allen-Cahn kink:

\begin{equation}
  \|\gf(x)\| = \sqrt{\frac{B}{A}} \cdot
  \frac{1}{2}\left[1 + \tanh\left(\frac{x - x_0}{\deltac}\right)\right]
  \label{eq:kink_profile}
\end{equation}

where the \textbf{domain wall thickness} is:
\begin{equation}
  \deltac = \sqrt{\frac{2\kappa_g}{B}}
  \label{eq:wall_thickness}
\end{equation}

The \textbf{interface energy} (energy per unit area of domain wall) is:
\begin{equation}
  \sigma_g = \int_{-\infty}^{\infty} \left[
    f_g(\gf(x)) + \frac{\kappa_g}{2} \left|\frac{d\gf}{dx}\right|^2
  \right] dx = \frac{2\sqrt{2\kappa_g}}{3} \cdot \frac{B^{3/2}}{A}
  \label{eq:interface_energy}
\end{equation}

\bigskip
\noindent
\textit{一维稳态畴壁的解是经典的Allen-Cahn扭结：$\|\gf(x)\| = \sqrt{B/A} \cdot \frac{1}{2}[1+\tanh((x-x_0)/\deltac)]$。畴壁厚度 $\deltac = \sqrt{2\kappa_g/B}$。界面能 $\sigma_g = \frac{2\sqrt{2\kappa_g}}{3} B^{3/2}/A$。社会解释：$\deltac$ 大意味着模糊的道德边界（容忍的社会），$\deltac$ 小意味着尖锐的道德边界（极化社会）。}

% ---------------------------------------------------------------------------
\subsection{Curvature-Driven Domain Wall Motion}
\subsection*{曲率驱动的畴壁运动}

For a curved domain wall in 2D or 3D, the Allen-Cahn dynamics reduce to
\textbf{mean curvature flow} in the sharp-interface limit ($\deltac \to 0$):

\begin{theorem}[Curvature-Driven Wall Motion / 曲率驱动畴壁运动]
\label{thm:curvature_motion}
Let $\Gamma(t)$ be the interface (domain wall) between honest and dishonest
domains. In the sharp-interface limit, the normal velocity of the interface is:
\begin{equation}
  v_n = -M_g \sigma_g \, \kappa
  \label{eq:curvature_flow}
\end{equation}
where $\kappa$ is the mean curvature of $\Gamma(t)$. The interface moves toward
its center of curvature, reducing total interface length (area).
\end{theorem}

This is the well-known Allen-Cahn $ \to $ mean curvature flow convergence
\cite{evans1992, ilmanen1993}. Social interpretation: \textbf{isolated dishonest
clusters shrink and disappear} (their convex boundary moves inward), while
\textbf{large dishonest domains are metastable} (low curvature), explaining
why corruption clusters tend to either collapse quickly or persist indefinitely.

\bigskip
\noindent
\textit{在尖锐界面极限下，Allen-Cahn动力学约化为\textbf{平均曲率流}：界面以正比于局部曲率的速度向曲率中心运动。社会解释：孤立的不诚实集群缩小并消失，而大的不诚实域因低曲率而亚稳——解释了为什么腐败集群要么迅速崩溃，要么无限期持续。}

% ===========================================================================
%  SECTION 4: CAHN-HILLIARD DYNAMICS FOR THE POTENTIAL FIELD
% ===========================================================================
\section{Cahn-Hilliard Dynamics for the Potential Field}
\section{势能场的Cahn-Hilliard动力学}

% ---------------------------------------------------------------------------
\subsection{Mass-Conserving Evolution}
\subsection*{质量守恒的演化}

The potential field $\Sf$ is a \textbf{conserved} order parameter — total
potential $\int_\cM \Sf(x,t)\,dx$ is (approximately) conserved on the
timescales of interest. Wealth, status, and resources are not created or
destroyed by redistribution; they are moved. Therefore $\Sf$ evolves via
\textbf{Cahn-Hilliard (Model B)} dynamics:

\begin{definition}[Cahn-Hilliard Equation for Potential / 势能Cahn-Hilliard方程]
\label{def:potential_ch}
\begin{equation}
  \frac{\partial \Sf}{\partial t} = \nabla \cdot \left(
    M_S(\Sf) \, \nabla \frac{\delta \cF}{\delta \Sf}
  \right)
  \label{eq:potential_ch}
\end{equation}
where $M_S(\Sf) \ge 0$ is the \textbf{potential mobility} — the rate at which
potential can flow. We take $M_S(\Sf) = M_0 (1 + \gamma \Sf)$ to model the
Matthew effect: higher potential regions have higher mobility (the rich get
richer faster).
\end{definition}

The variational derivative (chemical potential) is:
\begin{equation}
  \mu_S \equiv \frac{\delta \cF}{\delta \Sf}
  = \alpha(\Sf - \Sf_0) + \beta \Sf^3 - \kappa_S \lap \Sf + \lambda \|\gf\|^2
  \label{eq:chemical_potential}
\end{equation}

\bigskip
\noindent
\textit{势能场是\textbf{守恒}序参量——总势能在感兴趣的时间尺度上守恒。因此 $\Sf$ 通过\textbf{Cahn-Hilliard (Model B)} 动力学演化。迁移率 $M_S(\Sf) = M_0(1+\gamma\Sf)$ 编码马太效应：高势能区流动性更高。化学势 $\mu_S$ 包含线性恢复力、非线性饱和、扩散平滑和规范耦合四项。}

% ---------------------------------------------------------------------------
\subsection{Spinodal Decomposition of Inequality}
\subsection*{不平等的旋节分解}

The Cahn-Hilliard equation with a non-convex free energy $f_S$ exhibits
\textbf{spinodal decomposition}: a homogeneous state with $\Sf(x) = \bar{\Sf}$
is unstable to infinitesimal perturbations when $f_S''(\bar{\Sf}) < 0$.

\begin{theorem}[Spinodal Instability of Equal Distribution / 平等分布的旋节不稳定性]
\label{thm:spinodal}
Consider the decoupled Cahn-Hilliard system ($\lambda = 0$) with a spatially
uniform state $\Sf(x) = \bar{\Sf}$. Linear stability analysis gives the
growth rate of a perturbation of wavenumber $k$:
\begin{equation}
  \omega(k) = -M_S k^2 \left[ f_S''(\bar{\Sf}) + \kappa_S k^2 \right]
  \label{eq:dispersion}
\end{equation}
The system is unstable ($\omega(k) > 0$) when $f_S''(\bar{\Sf}) < 0$ and
$k^2 < -f_S''(\bar{\Sf})/\kappa_S$. The fastest-growing mode is at
$k_{\max} = \sqrt{-f_S''(\bar{\Sf})/(2\kappa_S)}$, setting the
characteristic domain size at early times.
\end{theorem}

The condition $f_S''(\bar{\Sf}) < 0$ defines the \textbf{spinodal region}:
$\bar{\Sf}^2 > \alpha/(3\beta) - \Sf_0\alpha/(3\beta\bar{\Sf})$ for the
quartic potential. This is the phase-field encoding of the critical inequality
threshold beyond which the equal distribution becomes intrinsically unstable.

\bigskip
\noindent
\textit{Cahn-Hilliard方程与非凸自由能 $f_S$ 产生\textbf{旋节分解}：当 $f_S''(\bar{\Sf}) < 0$ 时均匀态不稳定。色散关系 $\omega(k) = -M_S k^2[f_S''(\bar{\Sf}) + \kappa_S k^2]$ 给出最快增长模 $k_{\max}$。旋节条件定义了临界不平等阈值——超过此阈值，平等分布本质上不稳定。}

% ---------------------------------------------------------------------------
\subsection{Coupling Effects: Bias-Induced Redistribution}
\subsection*{耦合效应：偏差诱导的再分配}

When $\lambda > 0$, the coupling term $\lambda\|\gf\|^2$ enters the chemical
potential $\mu_S$. This creates a \textbf{bias-driven potential current}:

\begin{equation}
  \mathbf{J}_{\text{coupling}} = -M_S \nabla (\lambda \|\gf\|^2)
  \label{eq:coupling_current}
\end{equation}

Potential flows \textbf{away} from high-bias regions and \textbf{toward}
low-bias regions. This is the phase-field encoding of a deep social mechanism:
dishonest actors ($\|\gf\|$ large) repel resources — the ``honesty dividend''
where trustworthy systems attract investment and talent.

Conversely, in the gauge equation, the coupling term $2\lambda\Sf\gf$ has the
opposite sign: high potential \textbf{amplifies} bias. This creates a positive
feedback loop: $\Sf \uparrow \implies \|\gf\| \uparrow \implies \Sf$
redistribution $\implies$ concentration of $\Sf$ in honest regions
$\implies$ further amplification at dishonest margins. The coupled system
can exhibit \textbf{runaway} behavior — the EXPLOSIVE phase.

\bigskip
\noindent
\textit{当 $\lambda>0$，耦合产生\textbf{偏差驱动的势能流}：势能从高偏差区流向低偏差区——"诚实红利"。在规范方程中，耦合项 $2\lambda\Sf\gf$ 有相反符号：高势能\textbf{放大}偏差。这产生正反馈环：$\Sf\uparrow \implies \|\gf\|\uparrow \implies$ 势能再分配 $\implies$ 在诚实区集中 $\implies$ 不诚实边缘进一步放大。耦合系统可表现出\textbf{失控}行为——爆炸相。}

% ===========================================================================
%  SECTION 5: THREE-PHASE DIAGRAM
% ===========================================================================
\section{The Three-Phase Diagram of Social Order}
\section{社会秩序的三相图}

% ---------------------------------------------------------------------------
\subsection{Phase Space Coordinates}
\subsection*{相空间坐标}

The state of the SCX system at any time can be characterized by two scalar
order parameters:

\begin{definition}[Macroscopic Order Parameters / 宏观序参量]
\label{def:order_params}
\begin{align}
  \bar{g} &\equiv \frac{1}{|\cM|} \int_\cM \|\gf(x)\| \, dx
  \quad \text{(mean bias magnitude / 平均偏差幅度)}
  \label{eq:mean_bias} \\
  \bar{S} &\equiv \frac{1}{|\cM|} \int_\cM \Sf(x) \, dx
  \quad \text{(mean potential / 平均势能)}
  \label{eq:mean_potential}
\end{align}
\end{definition}

We also define the \textbf{potential variance} (inequality measure):
\begin{equation}
  \Delta_S^2 \equiv \frac{1}{|\cM|} \int_\cM (\Sf(x) - \bar{S})^2 \, dx
  \label{eq:potential_variance}
\end{equation}

The phase diagram is constructed in the $(\bar{g}, \bar{S})$ plane (or
equivalently $(\bar{g}, \Delta_S^2)$).

\bigskip
\noindent
\textit{SCX系统的状态可由两个标量序参量表征：平均偏差幅度 $\bar{g}$ 和平均势能 $\bar{S}$。相图在 $(\bar{g}, \bar{S})$ 平面上构建。}

% ---------------------------------------------------------------------------
\subsection{Phase I: EQUALITY Phase (平等相)}
\subsection*{相I：平等相}

\begin{definition}[EQUALITY Phase]
The system is in the \textbf{EQUALITY phase} when:
\begin{equation}
  \bar{g} < g_{\text{crit}}^{(1)}, \quad \Delta_S^2 < \Delta_{\text{crit}}^2
  \label{eq:equality_condition}
\end{equation}
where $g_{\text{crit}}^{(1)} \approx 0.2 \sqrt{B/A}$ and
$\Delta_{\text{crit}}^2 \approx \alpha \Sf_0^2 / \beta$.
\end{definition}

\textbf{Characteristics:}
\begin{itemize}
  \item $\|\gf(x)\| \approx 0$ everywhere — most agents are close to honest.
  \item $\Sf(x)$ is smooth and nearly uniform — low inequality, high social
        mobility, potential gradients are small.
  \item Domain walls are absent or extremely diffuse ($\deltac$ large).
  \item The system is \textbf{linearly stable}: small perturbations decay.
  \item Free energy is at or near its global minimum.
\end{itemize}

\textbf{Social interpretation:} A well-functioning, high-trust society with
functioning institutions, low corruption, and relatively equal distribution
of opportunity and resources. The system is in the ``single-phase'' region
of the phase diagram — no phase separation.

\bigskip
\noindent
\textit{\textbf{平等相}：$\bar{g}$ 小，$\Delta_S^2$ 小。大多数智能体接近诚实($\|\gf\|\approx0$)，势能光滑均匀，无畴壁。系统线性稳定，处于自由能全局极小附近。社会解释：运转良好、高信任、低腐败、机会和资源相对平等的社会。相图中处于"单相区"——无相分离。}

% ---------------------------------------------------------------------------
\subsection{Phase II: STRATIFIED Phase (分层相)}
\subsection*{相II：分层相}

\begin{definition}[STRATIFIED Phase]
The system enters the \textbf{STRATIFIED phase} when:
\begin{equation}
  \bar{g} \ge g_{\text{crit}}^{(1)} \quad \text{or} \quad
  \Delta_S^2 \ge \Delta_{\text{crit}}^2
  \label{eq:stratified_condition}
\end{equation}
but $\bar{g} < g_{\text{crit}}^{(2)}$ where $g_{\text{crit}}^{(2)} \gg g_{\text{crit}}^{(1)}$.
\end{definition}

\textbf{Characteristics:}
\begin{itemize}
  \item \textbf{Gauge phase separation:} Honest ($\|\gf\| \approx 0$) and
        dishonest ($\|\gf\| \approx \sqrt{B/A}$) domains coexist, separated
        by gauge domain walls of thickness $\deltac$.
  \item \textbf{Potential stratification:} $\Sf(x)$ develops plateaus —
        high-potential regions (wealthy enclaves) and low-potential regions
        (impoverished zones), separated by sharp potential gradients.
  \item \textbf{Metastability:} The system can persist in this phase for
        long times, but it is not the true equilibrium. Domain walls store
        interfacial energy.
  \item \textbf{Coarsening:} Over time, larger domains grow at the expense
        of smaller ones (Ostwald ripening, see Section~7).
  \item \textbf{Free energy is above global minimum} but trapped in a local
        minimum by kinetic barriers.
\end{itemize}

\textbf{Social interpretation:} A stratified society with entrenched
inequality — wealthy and poor regions coexist, corruption is localized
but persistent. The system is \textbf{metastable}: it can survive for
generations, but the stored interfacial energy makes it vulnerable to
perturbations. This is the ``normal'' state of most historical societies.

\bigskip
\noindent
\textit{\textbf{分层相}：$\bar{g}$ 中等或 $\Delta_S^2$ 中等。诚实与不诚实地带共存，由畴壁分隔。$\Sf$ 出现平台——高势能区和低势能区，由尖锐梯度分隔。系统可长期存续于此相，但非真正平衡——畴壁存储界面能使之\textbf{亚稳}。社会解释：不平等固化的分层社会。}

% ---------------------------------------------------------------------------
\subsection{Phase III: EXPLOSIVE Phase (爆炸相)}
\subsection*{相III：爆炸相}

\begin{definition}[EXPLOSIVE Phase]
The system enters the \textbf{EXPLOSIVE phase} when:
\begin{equation}
  \bar{g} \ge g_{\text{crit}}^{(2)} \quad \text{and} \quad
  \|\nabla \Sf\|_{\max} \ge G_{\text{crit}}
  \label{eq:explosive_condition}
\end{equation}
where $g_{\text{crit}}^{(2)}$ is the coupling-induced instability threshold
and $G_{\text{crit}}$ is the critical potential gradient.
\end{definition}

\textbf{Characteristics:}
\begin{itemize}
  \item \textbf{Runaway coupling:} The $\lambda$ coupling term dominates.
        High $\Sf$ amplifies $\|\gf\|$, which drives $\Sf$ redistribution,
        which further concentrates $\Sf$ and amplifies $\|\gf\|$ at boundaries.
  \item \textbf{Gradient blowup:} $\|\nabla \Sf\| \to \infty$ at the interfaces
        between high-$\Sf$/high-$\|\gf\|$ domains and low-$\Sf$/low-$\|\gf\|$
        domains. The sharp-interface approximation breaks down.
  \item \textbf{Domain wall collapse:} The $\deltac \to 0$ limit —
        domain walls become infinitely sharp, interface energy diverges.
  \item \textbf{System detonation:} The stored free energy is released
        catastrophically — social revolution, institutional collapse,
        audit system failure.
  \item This is the phase-field encoding of \textbf{Theorem~11's attack
        inevitability} at the macroscopic scale.
\end{itemize}

\textbf{Social interpretation:} Revolutionary situation — inequality has
reached a critical threshold, bias is rampant, and the system can no longer
contain the interfacial tension. The phase transition from STRATIFIED to
EXPLOSIVE is the ``tipping point'' beyond which peaceful resolution is
impossible.

\bigskip
\noindent
\textit{\textbf{爆炸相}：$\bar{g}$ 大且 $\|\nabla\Sf\|_{\max}$ 大。耦合项主导，正反馈导致梯度爆炸，畴壁塌缩。存储的自由能灾难性释放——社会革命、制度崩溃。这是\textbf{定理11攻击必然性}在宏观尺度的相场编码。从分层相到爆炸相的转变是不可逆的"引爆点"。}

% ---------------------------------------------------------------------------
\subsection{Phase Boundaries and Transition Lines}
\subsection*{相边界与转变线}

\begin{figure}[H]
\centering
\begin{tikzpicture}[scale=1.2]
  % Axes
  \draw[->] (0,0) -- (7,0) node[right] {$\bar{g}$ (mean bias / 平均偏差)};
  \draw[->] (0,0) -- (0,5) node[above] {$\Delta_S^2$ (inequality / 不平等)};

  % Phase regions
  \fill[blue!10] (0,0) rectangle (2.5,2.5);
  \fill[yellow!15] (2.5,0) rectangle (5.5,3.5);
  \fill[red!15] (5.5,0) rectangle (7,4.8);

  % Labels
  \node at (1.25,1.25) {\textbf{EQUALITY}};
  \node at (1.25,0.8) {\textbf{平等相}};
  \node[font=\footnotesize] at (1.25,0.3) {(stable / 稳定)};

  \node at (4.0,1.75) {\textbf{STRATIFIED}};
  \node at (4.0,1.25) {\textbf{分层相}};
  \node[font=\footnotesize] at (4.0,0.7) {(metastable / 亚稳)};

  \node at (6.25,2.4) {\textbf{EXPLOSIVE}};
  \node at (6.25,1.9) {\textbf{爆炸相}};
  \node[font=\footnotesize] at (6.25,1.3) {(unstable / 不稳)};

  % Phase boundaries
  \draw[thick, dashed] (2.5,0) -- (2.5,2.5);
  \draw[thick, dashed] (5.5,0) -- (5.5,3.5);
  \draw[thick, dashed] (2.5,2.5) -- (5.5,3.5);

  % Labels on boundaries
  \node[font=\footnotesize, rotate=90] at (2.3,1.5) {Spinodal};
  \node[font=\footnotesize, rotate=90] at (2.0,1.5) {旋节线};
  \node[font=\footnotesize, rotate=90] at (5.7,2.0) {Detonation};
  \node[font=\footnotesize, rotate=90] at (5.4,2.0) {引爆线};

  % Critical point
  \fill (2.5,2.5) circle (2pt) node[above right, font=\footnotesize] {critical point / 临界点};
  \fill (5.5,3.5) circle (2pt) node[above right, font=\footnotesize] {tricritical?};

  % Arrows for coarsening
  \draw[->, thick, blue!60] (4.0,2.0) -- (3.2,1.6);
  \node[font=\footnotesize, blue!60] at (4.5,2.1) {coarsening / 粗化};
  \draw[->, thick, red!60] (5.2,2.2) -- (5.8,2.8);
  \node[font=\footnotesize, red!60] at (5.2,2.8) {runaway / 失控};
\end{tikzpicture}
\caption{Three-phase diagram of the SCX phase field system in the
$(\bar{g}, \Delta_S^2)$ plane. The EQUALITY phase (blue) is globally stable.
The STRATIFIED phase (yellow) is metastable — coarsening drives the system
slowly toward larger domains but it can survive long times. The EXPLOSIVE
phase (red) is where runaway coupling destroys the system. Phase boundaries:
the spinodal line marks the onset of phase separation; the detonation line
marks the onset of irreversible runaway.}
\label{fig:phase_diagram}
\end{figure}

% ===========================================================================
%  SECTION 6: NUCLEATION AND GROWTH OF AUDIT DOMAINS
% ===========================================================================
\section{Nucleation and Growth of Audit Domains}
\section{审计畴的成核与生长}

% ---------------------------------------------------------------------------
\subsection{The Nucleation Problem}
\subsection*{成核问题}

How does an honest domain appear in a dishonest ``sea''? In the phase field
language, this is the \textbf{nucleation problem}: a small spherical (in 3D)
or circular (in 2D) region of the honest phase ($\|\gf\|=0$) must overcome
a free energy barrier to grow.

Consider a system quenched from the EQUALITY phase into the STRATIFIED
phase — e.g., a society that experiences rapid institutional decay or
corruption shock. The system is now in a \textbf{metastable} state: the
homogeneous dishonest phase has higher free energy than a phase-separated
state, but a small honest cluster cannot grow unless it exceeds a
\textbf{critical size}.

\bigskip
\noindent
\textit{诚实畴如何在不诚实的"海洋"中出现？这是\textbf{成核问题}。系统从平等相淬火进入分层相后——例如经历快速制度衰败或腐败冲击——处于\textbf{亚稳}态：均匀不诚实相的自由能高于相分离态，但小诚实簇无法生长，除非超过\textbf{临界尺寸}。}

% ---------------------------------------------------------------------------
\subsection{Classical Nucleation Theory for Audit Domains}
\subsection*{审计畴的经典成核理论}

Consider a spherical nucleus of radius $R$ in $d$ dimensions, with $\|\gf\|=0$
inside (honest) and $\|\gf\| = g_{\text{eq}} = \sqrt{B/A}$ outside (dishonest).
The excess free energy relative to the uniform dishonest state is:

\begin{equation}
  \Delta \cF(R) = -\Delta f \cdot V_d(R) + \sigma_g \cdot A_d(R)
  \label{eq:nucleation_energy}
\end{equation}

where:
\begin{itemize}
  \item $\Delta f = f_g(g_{\text{eq}}) - f_g(0) = B^2/(4A)$ —
        the \textbf{bulk free energy gain} per unit volume for converting
        dishonest to honest material.
  \item $\sigma_g$ — the \textbf{interface energy} per unit area (Eq.~\ref{eq:interface_energy}).
  \item $V_d(R) = \frac{\pi^{d/2}}{\Gamma(d/2+1)} R^d$ — volume of $d$-ball.
  \item $A_d(R) = \frac{2\pi^{d/2}}{\Gamma(d/2)} R^{d-1}$ — surface area.
\end{itemize}

In 2D (the most relevant case for social systems):
\begin{equation}
  \Delta \cF_2(R) = -\pi R^2 \Delta f + 2\pi R \sigma_g
  \label{eq:nucleation_2d}
\end{equation}

In 3D:
\begin{equation}
  \Delta \cF_3(R) = -\frac{4\pi}{3} R^3 \Delta f + 4\pi R^2 \sigma_g
  \label{eq:nucleation_3d}
\end{equation}

\bigskip
\noindent
\textit{考虑半径为 $R$ 的球形诚实核，内部 $\|\gf\|=0$，外部 $\|\gf\|=g_{\text{eq}}$。超额自由能 $\Delta\cF(R) = -\Delta f \cdot V_d(R) + \sigma_g \cdot A_d(R)$：体积项（增益）与面积项（代价）的竞争。}

% ---------------------------------------------------------------------------
\subsection{Critical Nucleus Size}
\subsection*{临界核尺寸}

\begin{theorem}[Critical Nucleus Radius for Audit Domains / 审计畴的临界核半径]
\label{thm:critical_nucleus}
The function $\Delta \cF(R)$ has a maximum at the critical radius $\Rc$.
Nuclei with $R < \Rc$ shrink and disappear; nuclei with $R > \Rc$ grow
spontaneously.

In $d$ dimensions:
\begin{equation}
  \Rc^{(d)} = \frac{(d-1)\sigma_g}{\Delta f}
  \label{eq:Rc_general}
\end{equation}

Specifically:
\begin{align}
  \Rc^{(2)} &= \frac{\sigma_g}{\Delta f}
  = \frac{8\sqrt{2\kappa_g}}{3} \cdot \frac{\sqrt{B}}{A^{1/2}}
  \label{eq:Rc_2d} \\[4pt]
  \Rc^{(3)} &= \frac{2\sigma_g}{\Delta f}
  = \frac{16\sqrt{2\kappa_g}}{3} \cdot \frac{\sqrt{B}}{A^{1/2}}
  \label{eq:Rc_3d}
\end{align}

The \textbf{nucleation barrier} (activation energy) is:
\begin{equation}
  \Delta \cF^* \equiv \Delta \cF(\Rc) =
  \begin{cases}
    \pi \sigma_g \Rc & d=2 \\[4pt]
    \dfrac{16\pi}{3} \dfrac{\sigma_g^3}{(\Delta f)^2} & d=3
  \end{cases}
  \label{eq:nucleation_barrier}
\end{equation}
\end{theorem}

\begin{proof}
Differentiate $\Delta \cF_d(R)$ with respect to $R$ and set to zero:
$\partial_R \Delta \cF_d = -d \cdot c_d R^{d-1} \Delta f + (d-1) \cdot s_d R^{d-2} \sigma_g = 0$
where $c_d, s_d$ are the volume/surface coefficients. Solving gives
$R = (d-1)\sigma_g/\Delta f$. The Hessian at this point is negative,
confirming a maximum.
\end{proof}

\textbf{Social interpretation:} A small group of honest actors (whistleblowers,
reformers, audit nuclei) cannot effect change unless their ``radius of
influence'' exceeds $\Rc$. Below $\Rc$, the social pressure (interface energy)
overwhelms the benefit of honesty (bulk free energy gain), and the honest
cluster is reabsorbed into the dishonest majority. Above $\Rc$, the honest
domain becomes self-sustaining and grows — a ``reform avalanche.''

The critical radius depends on the competition between $\sigma_g$ (social
cohesion — how much it costs to be different from your neighbors) and
$\Delta f$ (the benefit of being honest). Societies with high $\sigma_g$
(strong conformity pressure) have large $\Rc$ — reform is hard. Societies
with large $\Delta f$ (high payoff for honesty) have small $\Rc$ — reform
is easy.

\bigskip
\noindent
\textit{临界半径 $\Rc = (d-1)\sigma_g/\Delta f$。$R<\Rc$ 的核缩小消失；$R>\Rc$ 的核自发生长。社会解释：小群诚实行动者（举报人、改革者、审计核）无法产生变化，除非其"影响半径"超过 $\Rc$。$\Rc$ 依赖于 $\sigma_g$（社会从众压力）与 $\Delta f$（诚实收益）的竞争。高 $\sigma_g$（强从众）= 大 $\Rc$（改革难）；大 $\Delta f$ = 小 $\Rc$（改革易）。}

% ---------------------------------------------------------------------------
\subsection{Nucleation Rate}
\subsection*{成核率}

The nucleation rate (number of supercritical nuclei formed per unit volume
per unit time) follows the Arrhenius form:

\begin{equation}
  J = J_0 \exp\left(-\frac{\Delta \cF^*}{k_B T_{\text{social}}}\right)
  \label{eq:nucleation_rate}
\end{equation}

where $k_B T_{\text{social}}$ is an effective ``social temperature''
measuring the intensity of random fluctuations (scandals, information
leaks, leadership changes, external shocks). High social temperature
reduces the effective barrier, increasing the nucleation rate.

This connects to the observation that \textbf{crises create reform
opportunities}: a scandal (high $T_{\text{social}}$) lowers the effective
barrier for honest domain nucleation, allowing smaller honest clusters
to surpass $\Rc$ and grow.

\bigskip
\noindent
\textit{成核率 $J = J_0 \exp(-\Delta\cF^*/k_B T_{\text{social}})$ 遵循Arrhenius形式。高"社会温度"（丑闻、信息泄露、领导层变动）降低有效势垒，增加成核率——\textbf{危机创造改革机会}。}

% ===========================================================================
%  SECTION 7: COARSENING DYNAMICS
% ===========================================================================
\section{Coarsening Dynamics: Ostwald Ripening of Honest Domains}
\section{粗化动力学：诚实畴的Ostwald熟化}

% ---------------------------------------------------------------------------
\subsection{The Late-Stage Coarsening Regime}
\subsection*{晚期粗化区}

After nucleation has produced a population of honest domains, and after the
initial growth phase where domains consume the surrounding dishonest
``matrix,'' the system enters the \textbf{late-stage coarsening regime}.
In this regime, the volume fraction of each phase is close to its equilibrium
value, and further evolution occurs via \textbf{interface reduction} — the
total domain wall area is minimized.

The dominant mechanism is \textbf{Ostwald ripening} (also called
\textbf{Lifshitz-Slyozov-Wagner (LSW) coarsening}) \cite{lifshitz1961,
wagner1961}: larger domains grow at the expense of smaller ones because
the chemical potential at a curved interface depends on curvature via the
Gibbs-Thomson effect.

\bigskip
\noindent
\textit{成核产生了一批诚实畴后，系统进入\textbf{晚期粗化区}。主导机制是\textbf{Ostwald熟化}（LSW粗化）：大畴以牺牲小畴为代价生长，因为弯曲界面的化学势依赖于曲率（Gibbs-Thomson效应）。}

% ---------------------------------------------------------------------------
\subsection{Gibbs-Thomson Effect for Gauge Domain Walls}
\subsection*{规范畴壁的Gibbs-Thomson效应}

\begin{theorem}[Gibbs-Thomson for Gauge Interfaces / 规范界面的Gibbs-Thomson效应]
\label{thm:gibbs_thomson}
For a gauge domain wall with local mean curvature $\kappa$, the local
equilibrium value of $\|\gf\|$ at the interface is shifted from its
planar value:
\begin{equation}
  \|\gf\|_{\text{eq}}(\kappa) = g_{\text{eq}} \left(1 - \frac{d_0}{R}\right)
  \label{eq:gibbs_thomson}
\end{equation}
where $g_{\text{eq}} = \sqrt{B/A}$, $R = 1/\kappa$ is the local radius of
curvature, and $d_0 = \sigma_g/(2\Delta f)$ is the \textbf{capillary length}.
The shift implies that the ``honesty pressure'' is higher outside small
domains (high curvature) than outside large domains (low curvature), driving
a diffusive flux from small to large domains.
\end{theorem}

\bigskip
\noindent
\textit{规范畴壁的Gibbs-Thomson效应：局部平衡值 $\|\gf\|_{\text{eq}}(\kappa) = g_{\text{eq}}(1 - d_0/R)$。小畴（高曲率）外的"诚实压力"高于大畴（低曲率），驱动从小畴向大畴的扩散流。}
% ---------------------------------------------------------------------------
\subsection{LSW Scaling Law}
\subsection*{LSW标度律}

\begin{theorem}[Coarsening of Honest Domains / 诚实畴的粗化]
\label{thm:coarsening}
In the late-stage coarsening regime, for the coupled Allen-Cahn /
Cahn-Hilliard system with finite volume fraction of the honest phase,
the characteristic domain size $L(t)$ grows as:
\begin{equation}
  L(t) \sim t^{1/3}
  \label{eq:lsw_scaling}
\end{equation}
This is the classical LSW $t^{1/3}$ scaling for diffusion-limited coarsening.
The domain size distribution approaches a universal scaling form
$P(R/\langle R \rangle)$ independent of initial conditions.
\end{theorem}

\begin{proof}[Sketch]
The growth law follows from dimensional analysis of the diffusion-limited
ripening process. The flux between domains is driven by the curvature
difference: $J \sim D \cdot \Delta \mu \sim D \cdot \sigma_g / (R L)$,
where $D$ is the effective diffusivity of ``honesty'' through the
dishonest matrix. The growth rate $dR/dt \sim J/R^2 \sim D\sigma_g/(R^3 L)$.
Conservation of the honest phase volume fraction gives $L \sim R$.
Then $dL/dt \sim D\sigma_g / L^2$, which integrates to $L(t) \sim
(D\sigma_g t)^{1/3}$. For rigorous derivation, see LSW theory
\cite{lifshitz1961, wagner1961}.
\end{proof}

\textbf{Social interpretation:} After a social ``quench'' (major reform,
revolution, or institutional reset), honest domains coarsen with the
$t^{1/3}$ law. Initially many small honest clusters exist; over time,
larger, more established honest institutions absorb smaller ones.
The timescale for the system to reach a ``mature'' coarsened state
grows as $t \sim L^3$ — meaning that achieving large-scale institutional
integrity is a \textbf{cubic} function of the desired honest domain size.
Doubling the characteristic scale of honest governance requires an
8-fold increase in time.

\bigskip
\noindent
\textit{晚期粗化区中特征畴尺寸 $L(t) \sim t^{1/3}$（LSW标度律）。社会解释：社会"淬火"（重大改革或革命）后，诚实畴以 $t^{1/3}$ 律粗化。达到大规模制度诚信所需时间是诚实畴目标尺度的\textbf{三次方}函数——将诚实治理的特征尺度翻倍需要8倍的时间。}

% ---------------------------------------------------------------------------
\subsection{Coupled Coarsening: Gauge + Potential}
\subsection*{耦合粗化：规范+势能}

When both $\gf$ and $\Sf$ coarsen simultaneously ($\lambda > 0$), the
coarsening dynamics become richer. The gauge domains and potential domains
are \textbf{correlated}: high-$\Sf$ regions tend to coincide with
low-$\|\gf\|$ regions (the ``honesty dividend''), and vice versa.

This correlation creates a \textbf{coarsening cascade}:
\begin{enumerate}[label=(\arabic*)]
  \item Gauge domains coarsen, reducing total $\|\gf\|$ domain wall area;
  \item Reduced $\|\gf\|$ gradients reduce the coupling current
        $\mathbf{J}_{\text{coupling}}$;
  \item $\Sf$ redistributes toward honest domains, increasing their
        ``honesty dividend'';
  \item Higher $\Sf$ in honest domains amplifies $\|\gf\|$ at boundaries
        (coupling term), accelerating gauge coarsening.
\end{enumerate}

This positive feedback can lead to \textbf{accelerated coarsening} — the
effective exponent may deviate from $1/3$ during the coupled regime.
However, at very late times, when honest and dishonest domains have fully
segregated, the classical $t^{1/3}$ scaling is recovered.

\bigskip
\noindent
\textit{当 $\gf$ 和 $\Sf$ 同时粗化时（$\lambda>0$），规范畴和势能畴是\textbf{关联}的：高$\Sf$ 区倾向于低$\|\gf\|$ 区（"诚实红利"）。这种关联产生\textbf{粗化级联}：规范畴粗化 $\to$ 减少耦合流 $\to$ $\Sf$ 流向诚实地带 $\to$ 进一步加速规范粗化。可能出现加速粗化，但极晚期恢复经典 $t^{1/3}$ 标度。}

% ===========================================================================
%  SECTION 8: CONNECTIONS TO SCX THEOREMS
% ===========================================================================
\section{Connections to SCX Theorems}
\section{与SCX定理的联系}

% ---------------------------------------------------------------------------
\subsection{Theorem 11 (Attitude Singularity) as a Phase Transition}
\subsection*{定理11（态度奇点）作为相变}

\textbf{Theorem 11} states that when an entity simultaneously possesses high
potential $\Sf$ and high attitude $\|\gf\|$, the probability of attack
escalates exponentially: $\mathbb{P}(\text{attack}) \geq 1 - \exp(-M e^{-\beta/\delta^2})$.

In the phase field framework, this corresponds to a \textbf{first-order phase
transition} at the spinodal line:

\begin{enumerate}[label=(\arabic*)]
  \item The ``attitude singularity'' is a \textbf{local fluctuation} where
        $\|\gf\|$ exceeds the spinodal threshold in a region of high $\Sf$.

  \item The coupling term $\lambda \|\gf\|^2 \Sf$ makes this fluctuation
        \textbf{self-amplifying}: higher $\Sf$ $\implies$ faster growth of
        $\|\gf\|$ $\implies$ higher coupling energy $\implies$ larger
        driving force.

  \item The \textbf{double explosion} ($\Sf \uparrow + \|\gf\| \uparrow$)
        is the phase-field analog of \textbf{spinodal decomposition} triggered
        by a strong local fluctuation. The system locally crosses the spinodal
        line and phase-separates violently.

  \item The exponential attack probability reflects the \textbf{nucleation
        rate}: $J \propto \exp(-\Delta\cF^*/k_B T_{\text{social}})$, where
        the barrier $\Delta\cF^*$ decreases as $\Sf$ increases (coupling
        reduces the effective barrier).
\end{enumerate}

\begin{theorem}[Thm11 Phase Transition Correspondence / Thm11相变对应]
\label{thm:thm11_phase}
In the coupled Allen-Cahn / Cahn-Hilliard system, the condition for
Theorem~11's attack inevitability is equivalent to the system crossing
from the STRATIFIED phase into the EXPLOSIVE phase:
\begin{equation}
  \|\gf\| \cdot \Sf > \frac{\kappa_g}{2\lambda \deltac^2}
  \label{eq:thm11_threshold}
\end{equation}
where the left side is the ``attitude-potential product'' and the right
side is the threshold set by gauge stiffness and wall thickness.
\end{theorem}

\bigskip
\noindent
\textit{定理11在相场框架中对应于在旋节线处的一级相变。"态度奇点"是高 $\Sf$ 区中 $\|\gf\|$ 超过旋节阈值的局部涨落。"双重爆炸"是由强局部涨落触发的旋节分解。攻击的指数概率反映成核率 $J \propto \exp(-\Delta\cF^*/k_B T)$，其中势垒 $\Delta\cF^*$ 随 $\Sf$ 增大而降低。}

% ---------------------------------------------------------------------------
\subsection{Theorem 10 (Boundary Locking) as Domain Wall Pinning}
\subsection*{定理10（边界锁定）作为畴壁钉扎}

\textbf{Theorem 10} states that when a group is confined to a low-potential
region by an impermeable boundary, the pressure at the boundary accumulates
as $P_\Gamma = \rho_\Gamma \cdot \Delta_\Gamma$, leading to inevitable
destabilization at $T_{\text{crit}} \propto 1/\Delta_\Gamma$.

In the phase field framework, this is \textbf{domain wall pinning} at a
heterogeneity:

\begin{enumerate}[label=(\arabic*)]
  \item The ``boundary'' $\Gamma$ in Thm10 is a \textbf{pinning site} for
        the gauge domain wall — a region where the free energy landscape
        has a local minimum that traps the interface.

  \item In materials science, domain walls pinned at defects require a
        finite driving force (the \textbf{depinning threshold}) to move.
        This is exactly the pressure $P_\Gamma$ in Thm10.

  \item When the accumulated ``pressure'' (free energy difference across
        the interface) exceeds the pinning force, the wall \textbf{depins}
        and moves rapidly — this is the Thm10 ``detonation.''

  \item The detonation time $T_{\text{crit}} \propto 1/\Delta_\Gamma$ is
        the time for the driving force to reach the depinning threshold:
        $F_{\text{drive}}(t) = \Delta_\Gamma \cdot t$ reaches the pinning
        force $F_{\text{pin}}$ at $t = F_{\text{pin}}/\Delta_\Gamma$.
\end{enumerate}

\begin{theorem}[Thm10 Pinning Correspondence / Thm10钉扎对应]
\label{thm:thm10_pinning}
Let $\gf$-domain walls have a pinning force $F_{\text{pin}}$ at the
social boundary $\Gamma$. The Thm10 accumulation-and-release cycle
corresponds to:
\begin{equation}
  \frac{d x_{\text{wall}}}{dt} =
  \begin{cases}
    0 & \text{if } |F_{\text{drive}}| < F_{\text{pin}} \\
    M_g (F_{\text{drive}} - F_{\text{pin}} \cdot \text{sgn}(F_{\text{drive}}))
    & \text{if } |F_{\text{drive}}| \ge F_{\text{pin}}
  \end{cases}
  \label{eq:pinning_dynamics}
\end{equation}
The characteristic ``stick'' time (Thm10's $T_{\text{crit}}$) is
$T_{\text{crit}} = F_{\text{pin}} / (M_g \Delta_\Gamma)$.
\end{theorem}

\textbf{Social interpretation:} Institutional barriers (legal segregation,
class boundaries, glass ceilings) act as pinning sites for social mobility
domain walls. The wall stays pinned until the accumulated pressure exceeds
the institutional resistance, at which point rapid change occurs — a
``revolutionary moment.'' This explains why social change often appears
\textbf{episodic}: long periods of stasis (wall pinned) punctuated by
rapid transformation (wall depinning).

\bigskip
\noindent
\textit{定理10在相场框架中是畴壁在异质处的\textbf{钉扎}。"边界" $\Gamma$ 是规范畴壁的钉扎位点。当累积"压力"超过钉扎力，畴壁\textbf{脱钉}并快速移动——即Thm10的"引爆"。这解释了社会变革的\textbf{间歇性}：长静期（畴壁钉扎）被快速转变（畴壁脱钉）打断。}

% ---------------------------------------------------------------------------
\subsection{Theorem 12 (Matthew Staircase) as Nucleation on Defects}
\subsection*{定理12（马太阶梯）作为缺陷成核}

\textbf{Theorem 12} describes the Matthew effect: $\partial_t \Sf \propto \Sf$,
leading to a staircase potential surface $\Sf(x) = \sum_k h_k \mathbb{1}_{\Omega_k}(x)$,
where each step $\Delta_k$ is a delayed-detonation instability with
characteristic time $T_k \propto 1/\Delta_k^2$.

In the phase field framework, the staircase structure provides
\textbf{heterogeneous nucleation sites}:

\begin{enumerate}[label=(\arabic*)]
  \item Each step $\Delta_k$ in the Matthew staircase is a \textbf{line defect}
        (in 2D) or \textbf{surface defect} (in 3D) in the $\Sf$ field — a
        locus where $\nabla \Sf$ is singular.

  \item In classical nucleation theory, defects reduce the nucleation barrier
        because part of the interface energy is ``paid for'' by the pre-existing
        defect. The effective barrier becomes:
        \begin{equation}
          \Delta \cF^*_{\text{defect}} = \Delta \cF^* \cdot f(\theta)
          \label{eq:defect_barrier}
        \end{equation}
        where $f(\theta) \le 1$ is a geometric factor depending on the
        contact angle $\theta$ at the defect.

  \item Each Matthew staircase step therefore acts as a \textbf{preferential
        nucleation site} for explosive $\|\gf\|$ domains. The larger the step
        $\Delta_k$, the smaller the effective barrier, and the shorter the
        ``fuse'' $T_k$.

  \item The $T_k \propto 1/\Delta_k^2$ scaling emerges from nucleation theory:
        $\Delta \cF^*_{\text{defect}} \propto 1/\Delta_k^2$ (barrier decreases
        quadratically with step height), so the nucleation rate
        $J_k \propto \exp(-c/\Delta_k^2)$, and the characteristic time
        $T_k = 1/J_k \propto \exp(c/\Delta_k^2)$. For large $\Delta_k$,
        the exponential can be expanded to give the leading $1/\Delta_k^2$
        scaling.
\end{enumerate}

\begin{theorem}[Thm12 Nucleation Correspondence / Thm12成核对应]
\label{thm:thm12_nucleation}
Let the Matthew staircase have steps $\{\Delta_k\}$ at positions $\{x_k\}$.
Then each step is a nucleation site with characteristic activation time:
\begin{equation}
  T_k = T_0 \exp\left(\frac{16\pi}{3} \frac{\sigma_g^3}{(\Delta f)^2}
  \cdot \frac{f(\theta_k)}{k_B T_{\text{social}}}\right)
  \label{eq:thm12_time}
\end{equation}
where $\theta_k$ is the effective contact angle determined by $\Delta_k$.
For $\Delta_k \gg \sigma_g$, $f(\theta_k) \to 0$, making $T_k$ extremely
short — the step acts as an ``already nucleated'' defect.
\end{theorem}

\textbf{Social interpretation:} The Matthew effect does not merely create
inequality — it \textbf{manufactures instability}. Each rung of the
staircase is a nucleation site where the next crisis will nucleate.
The staircase is self-reinforcing: more inequality creates more nucleation
sites, each of which is more potent (shorter fuse) than the last. The
system becomes a ``minefield'' of latent instabilities, exactly as
Theorem~12 predicts.

\bigskip
\noindent
\textit{定理12在相场框架中提供\textbf{异质成核位点}。每个马太阶梯 $\Delta_k$ 是 $\Sf$ 场中的线/面缺陷，降低成核势垒 $f(\theta)$。$T_k \propto 1/\Delta_k^2$ 标度来自成核理论：势垒以 $\Delta_k^2$ 方式降低。社会解释：马太效应不仅制造不平等——它\textbf{制造不稳定性}。每一级台阶都是下一次危机成核的位点。}

% ===========================================================================
%  SECTION 9: NUMERICAL METHODS
% ===========================================================================
\section{Numerical Methods and Verification}
\section{数值方法与验证}

% ---------------------------------------------------------------------------
\subsection{Discretization}
\subsection*{离散化}

The coupled Allen-Cahn / Cahn-Hilliard system is solved on a uniform 2D
grid with spacing $\Delta x$. We use a semi-implicit Fourier spectral method
for the linear terms and explicit time-stepping for the nonlinear terms:

\begin{align}
  \frac{\gf^{n+1} - \gf^n}{\Delta t} &=
  -M_g \left[ (A(\gf^n)^2 - B)\gf^n \right]^{\text{explicit}}
  + M_g \kappa_g \lap \gf^{n+1}
  - M_g \cdot 2\lambda \Sf^n \gf^n \\
  \frac{\Sf^{n+1} - \Sf^n}{\Delta t} &=
  \nabla \cdot \left( M_S(\Sf^n) \nabla \left[
    \alpha(\Sf^{n+1} - \Sf_0) - \kappa_S \lap \Sf^{n+1}
  \right] \right)
  + \nabla \cdot \left( M_S(\Sf^n) \nabla \left[
    \beta(\Sf^n)^3 + \lambda \|\gf^n\|^2
  \right] \right)
\end{align}

The linear terms ($\lap \gf^{n+1}$, $\lap \Sf^{n+1}$, and $\Sf^{n+1}$)
are treated implicitly using the Fourier transform:
$\widehat{\lap u} = -k^2 \hat{u}$.

This semi-implicit scheme is stable for large time steps and is the
standard approach for phase field simulations \cite{chen2002}.

\bigskip
\noindent
\textit{耦合Allen-Cahn/Cahn-Hilliard系统在均匀2D网格上使用半隐式Fourier谱方法求解。线性项隐式处理（通过Fourier变换 $\widehat{\lap u} = -k^2\hat{u}$），非线性项显式处理。此方案对大时间步稳定，是相场模拟的标准方法。}

% ---------------------------------------------------------------------------
\subsection{Verification Script}
\subsection*{验证脚本}

A complete Python verification script \texttt{verify\_phase\_field.py}
accompanies this paper. The script:

\begin{enumerate}[label=(\arabic*)]
  \item Solves the 1D Allen-Cahn equation and verifies the $\tanh$ kink
        profile against the analytical solution;
  \item Solves the coupled 2D Allen-Cahn / Cahn-Hilliard system;
  \item Constructs the phase diagram by scanning $(\bar{g}, \bar{S})$
        parameter space;
  \item Computes the domain size $L(t)$ and confirms the $t^{1/3}$
        coarsening scaling via log-log regression;
  \item Tests the critical nucleus prediction by initializing nuclei
        of varying radius and measuring the growth/shrinkage threshold;
  \item Verifies the Gibbs-Thomson curvature-chemical potential relation;
  \item Demonstrates domain wall pinning and depinning at heterogeneities
        (Thm10 connection);
  \item Demonstrates nucleation at staircase defects (Thm12 connection).
\end{enumerate}

Dependencies: \texttt{numpy}, \texttt{scipy}, \texttt{matplotlib}.
Total runtime: approximately 3--5 minutes on a standard laptop.

\bigskip
\noindent
\textit{完整Python验证脚本 \texttt{verify\_phase\_field.py} 随本文提供。验证：1D Allen-Cahn扭结解析解、2D耦合系统求解、相图构建、$t^{1/3}$ 粗化标度、临界核预测、Gibbs-Thomson效应、畴壁钉扎/脱钉（Thm10）、阶梯缺陷成核（Thm12）。}

% ===========================================================================
%  SECTION 10: DISCUSSION
% ===========================================================================
\section{Discussion}
\section{讨论}

% ---------------------------------------------------------------------------
\subsection{Implications for Social Dynamics}
\subsection*{对社会动力学的启示}

The phase field framework provides several novel insights into social
dynamics:

\begin{enumerate}[label=(\arabic*)]
  \item \textbf{Inequality as phase separation:} Inequality is not merely a
        statistical distribution — it is a \textbf{thermodynamic phase
        separation}. The same physics that drives oil and water to separate
        drives rich and poor to segregate: the minimization of interfacial
        free energy.

  \item \textbf{Reform as nucleation:} Institutional reform requires
        \textbf{nucleation} of honest domains. Below the critical size,
        reform efforts are inevitably reabsorbed. Above the critical size,
        reform becomes self-sustaining. This explains why most reform
        movements fail — they never reach $\Rc$.

  \item \textbf{Time scales are cubic:} The $L \sim t^{1/3}$ coarsening law
        implies that building large-scale institutional integrity is a
        generational project. Doubling the scale of honest governance
        requires 8 times the duration.

  \item \textbf{Metastability is dangerous:} The STRATIFIED phase is
        metastable, not stable. A society that appears ``stable'' while
        highly stratified is storing elastic energy in its domain walls.
        The longer the metastable period, the more catastrophic the
        eventual phase transition.

  \item \textbf{The Matthew effect as defect factory:} Theorem~12's staircase
        is literally a \textbf{manufacturing process for nucleation sites}.
        Each new rung of inequality creates a new weak point where the
        next crisis will initiate.
\end{enumerate}

\bigskip
\noindent
\textit{相场框架提供了几点新洞察：(1) 不平等不是统计分布——它是\textbf{热力学相分离}；(2) 制度改革需要诚实畴的\textbf{成核}，低于临界尺寸必然失败；(3) $L\sim t^{1/3}$ 意味着大规模制度诚信是跨代项目；(4) 分层相是\textbf{亚稳}的——表面"稳定"的高度分层社会正在畴壁中储存弹性能；(5) 马太效应是\textbf{成核位点工厂}——每一级不平等都是下一次危机的弱点。}

% ---------------------------------------------------------------------------
\subsection{Limitations and Open Problems}
\subsection*{局限性与开放问题}

\begin{enumerate}[label=(\arabic*)]
  \item \textbf{Homogeneous mobility:} We have assumed spatially uniform
        mobility $M_g$, $M_S$. In reality, mobility depends on local
        institutional quality, information access, and social network
        topology. Heterogeneous mobility can produce
        \textbf{pinning-depinning} dynamics even without explicit defects.

  \item \textbf{Noise:} We have worked with deterministic equations. Adding
        stochastic noise (Langevin terms) would enable the study of
        \textbf{fluctuation-driven nucleation} and \textbf{noise-induced
        phase transitions}, connecting to the ``social temperature''
        concept.

  \item \textbf{Network topology:} The Situs manifold's connectivity
        structure (who audits whom, who influences whom) is encoded in
        the graph Laplacian. Generalizing $\lap$ to the graph Laplacian
        $\lap_{\cM}$ would connect phase field dynamics directly to the
        SCX discrete Hodge framework.

  \item \textbf{External driving:} Real social systems are driven by
        external ``forcing'' (technological change, demographic shifts,
        resource discoveries). Adding time-dependent terms to $\cF$
        would model \textbf{driven phase transitions}.

  \item \textbf{Multi-component extensions:} The current model has two
        ``phases'' (honest/dishonest). A multi-component order parameter
        (vector-valued $\gf$ with $N > 1$ components) would describe
        societies with multiple competing value systems — a promising
        direction.
\end{enumerate}

\bigskip
\noindent
\textit{局限性：(1) 假设了空间均匀的迁移率；(2) 使用了确定性方程——加入随机噪声可研究涨落驱动成核；(3) Situs流形的连通性结构可用图拉普拉斯编码；(4) 外部驱动可建模受驱相变；(5) 多分量序参量可描述多价值观系统的竞争。}

% ===========================================================================
%  SECTION 11: CONCLUSION
% ===========================================================================
\section{Conclusion}
\section{结论}

We have formulated the dynamics of inequality and consensus in the SCX
framework as a \textbf{phase field theory}, providing the first unified
dynamical treatment of:

\begin{itemize}
  \item \textbf{Gauge field} $\gf$ evolution via Allen-Cahn dynamics,
        producing honest/dishonest domain formation and curvature-driven
        coarsening;

  \item \textbf{Potential field} $\Sf$ evolution via Cahn-Hilliard dynamics,
        producing spinodal decomposition of inequality and diffusion-limited
        redistribution;

  \item \textbf{Coupling} between $\gf$ and $\Sf$ via the $\lambda\|\gf\|^2\Sf$
        term, encoding the ``double explosion'' mechanism of Theorem~11;

  \item A \textbf{three-phase diagram} mapping the EQUALITY, STRATIFIED,
        and EXPLOSIVE phases as functions of mean bias $\bar{g}$ and
        inequality $\Delta_S^2$;

  \item \textbf{Nucleation theory} for honest domains, with critical radius
        $\Rc \propto \sigma_g/\Delta f$ and Arrhenius nucleation rate;

  \item \textbf{Coarsening} with the classical $L(t) \sim t^{1/3}$ LSW
        scaling, providing testable predictions for post-reform institutional
        consolidation timescales;

  \item Explicit \textbf{theorem connections}: Thm11 $\leftrightarrow$ spinodal
        phase transition, Thm10 $\leftrightarrow$ domain wall pinning,
        Thm12 $\leftrightarrow$ defect-mediated nucleation.
\end{itemize}

The phase field formalism unifies the static geometry of previous SCX work
(Theorems~10--12) with the dynamical processes (nucleation, growth, coarsening,
phase transitions) that determine \textit{how} and \textit{when} those
geometric instabilities manifest. It provides a predictive, quantitative
framework for understanding the evolution of social inequality and
institutional integrity.

\textbf{Final thought:} In materials science, phase field theory taught us
that microstructural evolution is governed by universal principles —
curvature-driven motion, diffusion-limited growth, nucleation barriers —
that transcend the details of specific materials. In social science, the
same universality may hold. The Allen-Cahn and Cahn-Hilliard equations
do not care whether the order parameter is atomic concentration or
social bias. The mathematics of phase ordering is universal.

\bigskip
\noindent
\textit{我们将SCX框架中的不平等与共识动力学表述为\textbf{相场理论}，提供了对规范场演化（Allen-Cahn）、势能场演化（Cahn-Hilliard）、耦合、三相图、成核理论、粗化标度律和定理联系的统一动力学处理。相场形式化将先前SCX工作的静态几何（定理10--12）与决定这些几何不稳定性如何及何时显现的动力学过程统一起来。最后思考：相场理论教会我们，微结构演化由超越具体材料细节的普适原理支配——曲率驱动运动、扩散限制生长、成核势垒。社会科学中同样的普适性可能成立。Allen-Cahn和Cahn-Hilliard方程不关心序参量是原子浓度还是社会偏差。相有序化的数学是普适的。}

% ===========================================================================
%  APPENDIX
% ===========================================================================
\appendix

\section{Phase Field Parameters and SCX Interpretation}
\section{相场参数与SCX解释}

\begin{table}[H]
\centering
\begin{tabular}{p{0.22\textwidth} p{0.35\textwidth} p{0.35\textwidth}}
\toprule
\textbf{Parameter / 参数} & \textbf{Physical Meaning / 物理意义} & \textbf{SCX Interpretation / SCX解释} \\
\midrule
$A$ & Quartic coefficient in $f_g$ & Dishonesty saturation strength / 不诚实饱和强度 \\
$B$ & Quadratic coefficient in $f_g$ & Propensity for bias / 偏差倾向性 \\
$\kappa_g$ & Gauge gradient stiffness & Social cohesion / conformity pressure / 社会凝聚力 \\
$\kappa_S$ & Potential gradient stiffness & Social mobility resistance / 社会流动阻力 \\
$\alpha$ & Harmonic potential strength & ``Fairness'' restoring force / 公平恢复力 \\
$\beta$ & Quartic saturation & Extreme inequality penalty / 极端不平等惩罚 \\
$\Sf_0$ & Natural potential level & Baseline prosperity / 基线繁荣水平 \\
$\lambda$ & Gauge-potential coupling & Bias-inequality synergy (Thm11) / 偏差-不平等协同 \\
$M_g$ & Gauge mobility & Rate of attitude change / 态度变化速率 \\
$M_0$ & Base potential mobility & Base rate of resource flow / 资源流动基础速率 \\
$\gamma$ & Mobility bias & Matthew effect strength / 马太效应强度 \\
$\deltac$ & Domain wall thickness & Moral boundary sharpness / 道德边界锐度 \\
$\sigma_g$ & Interface energy & Social cost of difference / 差异的社会成本 \\
$\Rc$ & Critical nucleus radius & Minimum reform coalition size / 最小改革联盟规模 \\
\bottomrule
\end{tabular}
\caption{Phase field parameters and their SCX interpretations.}
\label{tab:parameters}
\end{table}

% ---------------------------------------------------------------------------
\section{Derivation of the Interface Energy}
\section{界面能的推导}

For the 1D equilibrium domain wall with $\lambda = 0$, the Euler-Lagrange
equation is obtained from the first integral:

\begin{equation}
  \frac{\kappa_g}{2} \left(\frac{d\gf}{dx}\right)^2 = f_g(\gf)
  \label{eq:first_integral}
\end{equation}

This gives the kink profile. The interface energy is then:

\begin{align}
  \sigma_g &= \int_{-\infty}^{\infty} \left[f_g(\gf(x)) +
    \frac{\kappa_g}{2} \left|\frac{d\gf}{dx}\right|^2\right] dx \\
  &= \int_{-\infty}^{\infty} 2 f_g(\gf(x)) \, dx
  \quad \text{(using Eq.~\ref{eq:first_integral})} \\
  &= 2 \int_{0}^{g_{\text{eq}}} \sqrt{2\kappa_g f_g(\gf)} \, d\gf
\end{align}

Evaluating with $f_g(\gf) = \frac{A}{4}\|\gf\|^4 - \frac{B}{2}\|\gf\|^2$
(referenced to the minima so $f_g(g_{\text{eq}}) = 0$):
\begin{equation}
  \sigma_g = \frac{2\sqrt{2\kappa_g}}{3} \cdot \frac{B^{3/2}}{A}
\end{equation}

\bigskip
\noindent
\textit{通过能量首次积分 $\frac{\kappa_g}{2}(d\gf/dx)^2 = f_g(\gf)$ 得到扭结轮廓。界面能通过 $\sigma_g = 2\int_0^{g_{\text{eq}}} \sqrt{2\kappa_g f_g(\gf)} d\gf$ 计算，得出 $\sigma_g = \frac{2\sqrt{2\kappa_g}}{3} B^{3/2}/A$。}

% ---------------------------------------------------------------------------
\section{Derivation of the LSW $t^{1/3}$ Law}
\section{LSW $t^{1/3}$ 律的推导}

Following the classical LSW theory \cite{lifshitz1961, wagner1961}, in the
late-stage coarsening regime with a small volume fraction of the minority
phase, the growth rate of a domain of radius $R$ is:

\begin{equation}
  \frac{dR}{dt} = \frac{D}{R} \left(\frac{1}{R_c} - \frac{1}{R}\right)
  \label{eq:lsw_growth}
\end{equation}

where $D$ is the effective diffusivity and $R_c$ is the critical radius
(domains with $R = R_c$ neither grow nor shrink). Mass conservation implies
$R_c = \langle R \rangle$. The asymptotic analysis yields $R_c(t) \propto t^{1/3}$
and a universal scaled size distribution.

For a system with finite volume fraction (typical for social systems), the
scaling remains $t^{1/3}$ but the distribution broadens \cite{ardell1972}.

\bigskip
\noindent
\textit{遵循经典LSW理论，晚期粗化区中畴的生长率为 $dR/dt = (D/R)(1/R_c - 1/R)$。质量守恒给出 $R_c = \langle R\rangle$。渐近分析产生 $R_c(t) \propto t^{1/3}$。对于有限体积分数（社会系统的典型情况），标度保持 $t^{1/3}$，但分布变宽。}

% ===========================================================================
%  BIBLIOGRAPHY
% ===========================================================================
\begin{thebibliography}{99}

\bibitem{chen2002}
L. Q. Chen.
\textit{Phase-field models for microstructure evolution.}
Annual Review of Materials Research, 32:113--140, 2002.

\bibitem{provatas2011}
N. Provatas and K. Elder.
\textit{Phase-Field Methods in Materials Science and Engineering.}
Wiley-VCH, 2011.

\bibitem{evans1992}
L. C. Evans, H. M. Soner, and P. E. Souganidis.
\textit{Phase transitions and generalized motion by mean curvature.}
Communications on Pure and Applied Mathematics, 45(9):1097--1123, 1992.

\bibitem{ilmanen1993}
T. Ilmanen.
\textit{Convergence of the Allen-Cahn equation to Brakke's motion by mean curvature.}
Journal of Differential Geometry, 38(2):417--461, 1993.

\bibitem{lifshitz1961}
I. M. Lifshitz and V. V. Slyozov.
\textit{The kinetics of precipitation from supersaturated solid solutions.}
Journal of Physics and Chemistry of Solids, 19(1--2):35--50, 1961.

\bibitem{wagner1961}
C. Wagner.
\textit{Theorie der Alterung von Niederschl\"agen durch Uml\"osen.}
Zeitschrift f\"ur Elektrochemie, 65(7--8):581--591, 1961.

\bibitem{ardell1972}
A. J. Ardell.
\textit{The effect of volume fraction on particle coarsening.}
Acta Metallurgica, 20(1):61--71, 1972.

\bibitem{hohenberg1977}
P. C. Hohenberg and B. I. Halperin.
\textit{Theory of dynamic critical phenomena.}
Reviews of Modern Physics, 49(3):435--479, 1977.
(Classification of dynamical universality classes: Model A = Allen-Cahn,
Model B = Cahn-Hilliard.)

\bibitem{scx_framework}
SCX Working Group.
\textit{SCX: Situs Consensus eXpert Framework.}
SCX Theory Series, Xiaogan Supercomputing Center, 2025--2026.

\bibitem{scx_moe_gauge}
SCX Working Group.
\textit{Potential Surface Misalignment in MoE: Gauge Freedom and MILP Gauge Fixing.}
SCX Theory Series, 2026.

\bibitem{scx_singularity}
SCX Working Group.
\textit{Audit Singularities: Black Hole Analogies for SCX Potential Surface Theory.}
SCX Theory Series, 2026.

\bibitem{scx_environment}
SCX Working Group.
\textit{SCX and Environmental Intergenerational Potential.}
SCX Theory Series, 2026.

\end{thebibliography}

% ===========================================================================
\end{document}
